\documentclass[11pt,aspectratio=169]{beamer}
\usetheme{SimplePlus}

\usepackage[T1]{fontenc}
\usepackage[utf8]{inputenc}
\usepackage{amsmath,amssymb,mathtools}
\usepackage{booktabs}
\usepackage{array}
\usepackage{appendixnumberbeamer}
\usepackage{hyperref}

\newcommand{\E}{\mathbb{E}}
\newcommand{\Var}{\mathrm{Var}}
\newcommand{\Cov}{\mathrm{Cov}}
\newcommand{\indep}{\perp\!\!\!\perp}
\newcommand{\derivbutton}[1]{\vspace{0.4em}\hfill\hyperlink{app:#1}{\beamerbutton{Appendix details}}}
\newcommand{\backbutton}[1]{\vspace{0.4em}\hfill\hyperlink{main:#1}{\beamerbutton{Back to main slide}}}

\title{Applied Econometrics: Session 1}
\subtitle{Foundations of Causal Inference}
\author{PhD Intensive Course}
\institute{Lecture 1}
\date{\today}

\begin{document}

\begin{frame}
  \titlepage
\end{frame}

\begin{frame}{How These Slides Are Structured}
  \begin{itemize}
    \item This deck follows the lecture sequence: "Questions about Questions" and "The Experimental Ideal".
    \item Every core idea appears in three layers: question, intuition, and formula.
    \item The target is self-study: students should be able to revise directly from these slides.
    \item Detailed algebra is in appendix with back-and-forth buttons.
  \end{itemize}
\end{frame}

\begin{frame}{Roadmap for Today}
  \tableofcontents
\end{frame}

\section{Questions About Questions}

\begin{frame}{Why Start with Questions, Not Methods?}
  \begin{itemize}
    \item Good econometrics cannot rescue a weak research question.
    \item A coherent research agenda is the first identification device.
    \item Experimental way of thinking before choosing estimators.
  \end{itemize}
  \begin{block}{Core message}
    First define the causal question and ideal experiment; only then choose empirical strategy.
  \end{block}
\end{frame}

\begin{frame}{The Four Research FAQs}
  \begin{enumerate}
    \item What is the causal relationship of interest?
    \item What is the ideal experiment that answers it?
    \item What is the identification strategy in available data?
    \item What is the mode of statistical inference?
  \end{enumerate}
  \vspace{0.3em}
  \textbf{These are sequential, not optional.}
\end{frame}

\begin{frame}{FAQ 1: What Is the Causal Relationship?}
  Example: effect of schooling on wages.
  \begin{itemize}
    \item Descriptive statement: high-school and college graduates earn differently.
    \item Causal statement: what would happen to the same person's wage if schooling changed?
    \item Policy relevance: tuition subsidies, compulsory schooling laws, financial aid.
  \end{itemize}
  \begin{block}{Practical rule}
    If the question has no counterfactual interpretation, causal methods are not yet the right tool.
  \end{block}
\end{frame}

\begin{frame}{FAQ 2: What Is the Ideal Experiment?}
  \begin{itemize}
    \item Imagine an RCT with no budget and no practical constraints.
    \item This clarifies what must vary and what must stay fixed.
    \item Even hypothetical experiments improve question quality.
  \end{itemize}
  \textbf{Example:}
  \begin{itemize}
    \item To estimate schooling returns, ideal design randomly changes years of schooling.
    \item Real world cannot do this perfectly, but the thought experiment defines the target.
  \end{itemize}
\end{frame}

\begin{frame}{FAQ 2 Extension: FUQ (Fundamentally Unidentified Questions)}
  \hypertarget{main:fuq}{}
  Some questions cannot be answered by any feasible experiment.
  \vspace{0.3em}
  \textbf{Example:}
  effect of school starting age on first-grade scores.
  \begin{itemize}
    \item Older starters are older at testing (maturation effect).
    \item If you compare at same age, they have different time in school.
    \item Age and schooling exposure are mechanically linked in child samples.
  \end{itemize}
  \begin{block}{Lesson}
    Sometimes data cannot identify the target because the target itself is not separable.
  \end{block}
  \derivbutton{fuq-algebra}
\end{frame}

\begin{frame}{FAQ 3: What Is the Identification Strategy?}
  \begin{itemize}
    \item Identification strategy means: how observational data mimic an experiment.
    \item Typical sources: policy rules, institutional thresholds, lotteries, timing differences.
    \item Strategy must be described before estimation details.
  \end{itemize}
  \textbf{Guiding principle:} design first, mechanics second.
\end{frame}

\begin{frame}{FAQ 4: What Is the Mode of Inference?}
  \begin{itemize}
    \item What population does the sample represent?
    \item What is the sampling/assignment process?
    \item How should standard errors be computed (heteroskedasticity, clustering, design effects)?
  \end{itemize}
  \begin{block}{Why this matters}
    You can have a strong design but still get misleading precision if inference is wrong.
  \end{block}
\end{frame}

\begin{frame}{Workflow for Applied Work}
  \begin{enumerate}
    \item Write the causal estimand in words and symbols.
    \item Describe an ideal randomized trial.
    \item State what real-world variation approximates that trial.
    \item Explain why inference assumptions match the data structure.
  \end{enumerate}
  \vspace{0.3em}
  This sequence prevents most beginner mistakes.
\end{frame}

% \section{The Experimental Ideal}

\begin{frame}{The Experimental Ideal: Motivation}
  \begin{itemize}
    \item Random assignment is the benchmark for causal credibility.
    \item Nonexperimental comparisons often point in the wrong direction.
    \item We need to understand why selection bias appears and how experiments remove it.
  \end{itemize}
\end{frame}

\begin{frame}{Hospital Example}
  Question: Do hospitals make people healthier?
  \begin{itemize}
    \item Naive data comparison often shows hospitalized people in worse health.
    \item This does not prove hospital care is harmful.
    \item Sicker people are more likely to go to hospital in the first place.
  \end{itemize}
  \begin{block}{Interpretation}
    Observed differences combine treatment effects and pre-treatment differences.
  \end{block}
\end{frame}

\begin{frame}{Potential Outcomes Notation}
  \hypertarget{main:po}{}
  For each unit \(i\):
  \[
  Y_{1i} = \text{outcome if treated}, \qquad Y_{0i} = \text{outcome if untreated}.
  \]
  Treatment indicator: \(D_i \in \{0,1\}\).
  Observed outcome:
  \[
  Y_i = Y_{0i} + (Y_{1i}-Y_{0i})D_i.
  \]
  \begin{itemize}
    \item Individual causal effect: \(Y_{1i}-Y_{0i}\).
    \item We never observe both potential outcomes for the same unit at the same time.
  \end{itemize}
\end{frame}

\begin{frame}{Fundamental Problem of Causal Inference}
  \begin{itemize}
    \item For treated units, \(Y_{0i}\) is missing.
    \item For untreated units, \(Y_{1i}\) is missing.
    \item Therefore causal inference is a missing-counterfactual problem.
  \end{itemize}
  \begin{block}{Key consequence}
    We need assumptions or design features to reconstruct missing counterfactual means.
  \end{block}
\end{frame}

\begin{frame}{Selection Bias Decomposition}
  \hypertarget{main:decomp}{}
  \[
  \E[Y_i\mid D_i=1]-\E[Y_i\mid D_i=0]
  =
  \underbrace{\E[Y_{1i}-Y_{0i}\mid D_i=1]}_{\text{ATT}}
  +
  \underbrace{\E[Y_{0i}\mid D_i=1]-\E[Y_{0i}\mid D_i=0]}_{\text{selection bias}}.
  \]
  \begin{itemize}
    \item Left side is observed mean difference.
    \item First term is causal effect for treated units.
    \item Second term is non-causal baseline difference between groups.
  \end{itemize}
  \derivbutton{decomp-proof}
\end{frame}

\begin{frame}{Numeric Intuition for Selection Bias}
  Suppose observed difference is \(-0.70\) health points (higher = worse health).
  \begin{itemize}
    \item Treated group is sicker even without treatment: selection term is \(-1.00\).
    \item Then treatment effect on treated is:
    \[
    \text{ATT} = -0.70 - (-1.00) = +0.30.
    \]
  \end{itemize}
  \begin{block}{Interpretation}
    Naive comparison says treatment looks harmful; corrected effect is beneficial.
  \end{block}
\end{frame}

\begin{frame}{Why Random Assignment Solves the Selection Problem}
  \hypertarget{main:rand}{}
  Random assignment gives:
  \[
  D_i \indep (Y_{1i},Y_{0i}).
  \]
  Therefore:
  \[
  \E[Y_{0i}\mid D_i=1]=\E[Y_{0i}\mid D_i=0]
  \]
  so selection bias term is zero.
  \[
  \E[Y_i\mid D_i=1]-\E[Y_i\mid D_i=0]=\E[Y_{1i}-Y_{0i}].
  \]
  \begin{itemize}
    \item Experiments create comparability by design, not by modeling assumptions.
  \end{itemize}
  \derivbutton{rand-proof}
\end{frame}

\begin{frame}{What Experiments Do Not Automatically Solve}
  \begin{itemize}
    \item Attrition can destroy comparability if dropout differs by treatment status.
    \item Noncompliance weakens treatment contrast (assignment vs receipt).
    \item Spillovers can violate stable treatment definitions.
    \item Measurement problems can still bias results.
  \end{itemize}
  \textbf{Conclusion:} randomization is powerful, but implementation details still matter.
\end{frame}

\begin{frame}{Case Study: Tennessee STAR}
  \begin{itemize}
    \item Large education RCT on class size in early grades.
    \item Students assigned to small class, regular class, or regular class plus aide.
    \item Goal: estimate causal effect of class size on achievement.
  \end{itemize}
  Why influential:
  \begin{itemize}
    \item Design is transparent.
    \item Policy relevance is immediate.
    \item Results differ from many naive observational comparisons.
  \end{itemize}
\end{frame}

\begin{frame}{STAR Balance and Interpretation}
  Typical first check in any RCT: covariate balance across groups.
  \begin{itemize}
    \item Race, age, and free-lunch status should be similar if randomization worked.
    \item Class-size differences should be large across arms (first-stage check).
    \item Attrition differences need attention because they can induce bias later.
  \end{itemize}
  \begin{block}{Empirical lesson}
    Even in RCTs, credibility comes from both assignment and execution quality.
  \end{block}
\end{frame}

\begin{frame}{Experimental Effects from Means and Regressions}
  In a two-group experiment, effect estimate is difference in means:
  \[
  \hat\tau = \bar Y_{D=1}-\bar Y_{D=0}.
  \]
  This is numerically equivalent to OLS coefficient from
  \[
  Y_i = \alpha + \rho D_i + u_i.
  \]
  \begin{itemize}
    \item Here \(\hat\rho=\hat\tau\).
    \item Regression is a convenient language for experimental analysis.
  \end{itemize}
  \derivbutton{ols-dim}
\end{frame}

\begin{frame}{Regression Analysis of Experiments}
  \hypertarget{main:reg-exp}{}
  With constant treatment effect model:
  \[
  Y_i = \alpha + \rho D_i + \eta_i,
  \]
  where \(\rho\) is causal effect and \(\eta_i\) captures untreated heterogeneity.
  \vspace{0.3em}
  \begin{itemize}
    \item In randomized data, \(D_i\) is orthogonal to \(\eta_i\).
    \item OLS recovers causal \(\rho\) without omitted variable bias.
    \item Adding pre-treatment controls can improve precision.
  \end{itemize}
  \derivbutton{reg-causal}
\end{frame}

\begin{frame}{From Ideal Experiments to Quasi-Experiments}
  \begin{itemize}
    \item True RCTs are often expensive, slow, or infeasible.
    \item Strategy: search for natural variation that is "as good as random".
    \item Benchmark remains the notional randomized trial.
  \end{itemize}
  Example:
  \begin{itemize}
    \item Class-size caps create discontinuous class-size changes around enrollment thresholds.
    \item Units near threshold can be highly comparable.
  \end{itemize}
\end{frame}

\begin{frame}{Chapter 2 Practical Checklist}
  \begin{enumerate}
    \item Define treatment and outcome precisely.
    \item Write potential outcomes and target estimand.
    \item Decompose observed difference into effect plus selection.
    \item Ask if assignment is random or plausibly as-if random.
    \item Audit threats: attrition, noncompliance, spillovers, inference.
  \end{enumerate}
\end{frame}

\section{Synthesis for Lecture 1}

\begin{frame}{Integrated Summary of Chapters 1 and 2}
  \begin{itemize}
    \item Chapter 1 gives the research-design blueprint (four FAQs).
    \item Chapter 2 gives the formal causal engine (potential outcomes + selection problem + random assignment).
    \item Regression enters as a practical estimator for experimental contrasts.
    \item Credible empirical work combines clear question, design logic, and valid inference.
  \end{itemize}
\end{frame}

\begin{frame}{What Students Should Be Able to Do After This Lecture}
  \begin{itemize}
    \item Distinguish causal questions from descriptive questions.
    \item Write and interpret the selection-bias decomposition.
    \item Explain why random assignment identifies average treatment effects.
    \item Understand why OLS with a treatment dummy equals difference in means in experiments.
    \item Use the four FAQs to structure research proposals.
  \end{itemize}
\end{frame}

\appendix

\begin{frame}{Appendix Map}
  \begin{itemize}
    \item A1: Selection-bias decomposition proof.
    \item A2: Random assignment identification proof.
    \item A3: OLS equals difference in means derivation.
    \item A4: Constant-effect regression and causal interpretation.
    \item A5: FUQ school-starting-age algebra.
  \end{itemize}
\end{frame}

\begin{frame}{Appendix A1: Selection-Bias Decomposition Proof}
  \hypertarget{app:decomp-proof}{}
  Start with observed gap:
  \[
  \E[Y_i\mid D_i=1]-\E[Y_i\mid D_i=0].
  \]
  \[
  \E[Y_i\mid D_i=1]=\E[Y_{1i}\mid D_i=1], \quad
  \E[Y_i\mid D_i=0]=\E[Y_{0i}\mid D_i=0].
  \]
  Add and subtract \(\E[Y_{0i}\mid D_i=1]\):
  \[
  \E[Y_{1i}-Y_{0i}\mid D_i=1]
  +
  \left(\E[Y_{0i}\mid D_i=1]-\E[Y_{0i}\mid D_i=0]\right).
  \]
  \backbutton{decomp}
\end{frame}

\begin{frame}{Appendix A2: Why Random Assignment Identifies Causal Effects}
  \hypertarget{app:rand-proof}{}
  If assignment is random:
  \[
  D_i \indep (Y_{1i},Y_{0i}).
  \]
  Then
  \[
  \E[Y_{1i}\mid D_i=1]=\E[Y_{1i}],\qquad
  \E[Y_{0i}\mid D_i=0]=\E[Y_{0i}].
  \]
  Therefore observed difference simplifies to
  \[
  \E[Y_i\mid D_i=1]-\E[Y_i\mid D_i=0]
  =
  \E[Y_{1i}]-\E[Y_{0i}]
  =
  \E[Y_{1i}-Y_{0i}].
  \]
  Selection bias is eliminated because untreated potential outcomes are balanced by design.
  \backbutton{rand}
\end{frame}

\begin{frame}{Appendix A3: OLS with Binary Treatment Equals Difference in Means}
  \hypertarget{app:ols-dim}{}
  Regression:
  \[
  Y_i=\alpha+\rho D_i+u_i,\quad D_i\in\{0,1\}.
  \]
  Predicted mean for controls (\(D=0\)):
  \[
  \E[Y\mid D=0]=\alpha.
  \]
  Predicted mean for treated (\(D=1\)):
  \[
  \E[Y\mid D=1]=\alpha+\rho.
  \]
  Subtract:
  \[
  \rho = \E[Y\mid D=1]-\E[Y\mid D=0].
  \]
  Sample OLS counterpart gives
  \[
  \hat\rho = \bar Y_{D=1}-\bar Y_{D=0}.
  \]
  \backbutton{reg-exp}
\end{frame}

\begin{frame}{Appendix A4: Constant-Effect Model and Causal Interpretation}
  \hypertarget{app:reg-causal}{}
  Suppose
  \[
  Y_i=Y_{0i}+\rho D_i,
  \]
  where \(\rho\) is constant across units.
  Write
  \[
  Y_i=\alpha+\rho D_i+\eta_i,
  \]
  with
  \[
  \alpha=\E[Y_{0i}],\qquad \eta_i=Y_{0i}-\E[Y_{0i}].
  \]
  If \(D_i\indep Y_{0i}\), then \(\Cov(D_i,\eta_i)=0\), so OLS recovers \(\rho\).
  \vspace{0.3em}
  In observational data this orthogonality can fail, generating omitted-variable bias.
  \backbutton{reg-exp}
\end{frame}

\begin{frame}{Appendix A5: FUQ Algebra for School Starting Age}
  \hypertarget{app:fuq-algebra}{}
  Let
  \[
  \text{Current age} = \text{start age} + \text{time in school}.
  \]
  If outcome is first-grade test score, then for children still in school:
  \begin{itemize}
    \item Varying start age changes current age at test.
    \item Holding current age fixed changes time in school.
  \end{itemize}
  So direct effect of start age is not separable from maturation and schooling exposure in this setting.
  \begin{block}{Key point}
    Some targets are structurally unidentified in the available comparison frame.
  \end{block}
  \backbutton{fuq}
\end{frame}

\end{document}
